% 18-research.tex: future engineering and research.
\section{Future research and engineering}
\label{sec:research}

Everything in this section is marked \sFut{} unless stated otherwise. Items are ordered by how
directly they bear on the future environment of Section~\ref{sec:agents}.

\subsection{Presentation technologies that would remove the correlator}

The stable credential value in a presentation is the one property that separates Polaris from a
privacy-preserving proof of personhood. Five families of technique would address it, each with a
cost the repository would have to accept.

\begin{description}
  \item[Pairwise identifiers.] A different handle per relying party, derived from a holder secret
    and the party's identity, so that two parties cannot join their logs. Needs a holder-side secret
    and a rule for how the issuer's status artifacts bind to a derived handle.
  \item[Blinded or derived presentations.] The holder proves possession of a signature over
    attributes without revealing the signature itself, so a reused signature is not a correlator.
    The classical constructions rely on pairing-based assumptions that a quantum computer breaks;
    post-quantum equivalents are an open research area, and Polaris's post-quantum posture would
    have to be preserved.
  \item[One-time presentation tokens.] Single-use artifacts minted per presentation, at the cost of
    issuer contact or a pre-fetched batch, which trades against offline availability and issuer
    non-observation exactly as the status assertion does.
  \item[Anonymous credentials with selective disclosure.] Revealing that an attribute satisfies a
    predicate without revealing the attribute. Polaris's disclosure levels are recorded at the
    issuer; a presentation-layer selective disclosure would be new machinery, and the repository
    explicitly says it does not claim to be one.
  \item[Zero-knowledge uniqueness.] The scoped nullifier of Section~\ref{sec:agents}: one added
    public input to the existing circuit, per relying party and per epoch, never global.
\end{description}

\subsection{The holder-side key}

Every item above, the browser bridge, agent delegation, and on-device proving share one missing
piece: a key the holder controls, registered against the credential at enrolment. The current model
is issuer-centric on purpose, to keep key custody, rotation and recovery from becoming a per-person
problem, and the notarial signing and possession-based login follow from it. A holder key would be a
credential extension with its own recovery ceremony, and the recovery layer of
Section~\ref{sec:core} is the obvious place it would attach. The sibling-path witness optimisation,
already named in the roadmap, would let a holder prove membership against a national-scale tree
without reconstructing it, which on-device proving requires.

\subsection{Metis: learning the system, never the people}

The assessment the roadmap requires before any code would have to falsify the central hypothesis
(that a model can infer properties of the system without inferring the rights of people) against a
synthetic-nation ground truth, compare every model to a deterministic rule, and design the
constitutional invariant that would forbid the prohibited uses structurally, in the manner of
Athena's five checks. The candidate uses, infrastructure anomalies, capacity forecasts with
uncertainty, regression and release-risk detection, chaos learning, root-cause ranking over Athena's
dependency graph, and aggregate institutional anomalies, are all about Polaris and institutions,
never about citizens.

\subsection{Myrmex: stigmergic debugging as an outside observer}

The observers of Section~\ref{sec:senses} all live inside the repository and inherit its
assumptions. This paper proposes, under the provisional name Myrmex, a debugging system that would
live outside Polaris and attack it independently, inspired by how an ant colony coordinates without
a central plan.

\figfile{stigmergy}{The proposed colony. Small specialised agents traverse the code and the
architecture and leave machine-readable pheromones on a persistent evidence landscape: suspicion, a
trust-boundary anomaly, an unusual state transition, a cryptographic assumption, a missing edge, an
inconsistent claim, an area repeatedly associated with failure. Pheromones strengthen under
independent evidence and decay unsupported; agents follow reinforced trails. An orchestrator allocates
work and aggregates evidence and decides nothing. Confidence comes only from external evidence, never
from agreement between models. What survives reproduction enters Polaris's own loop as a test and an
invariant.}{fig:stigmergy}

The roles would be specialised: scouts that map surfaces; trust-boundary agents that look for a
claim crossing a boundary without its proof; cryptographic agents that test each verdict against
the two-witness rule and the canonical discipline; privacy agents that hunt for a person-level
handle or an unbounded read; state-transition agents that try to reach a state the schema should
refuse; concurrency agents that race the advisory locks; protocol agents that mutate every signed
artifact the way the fuzzer does but with intent; and refutation agents whose only job is to
disprove what the others suspect. The queen, or orchestrator, allocates work and merges evidence
and is deliberately not an omniscient authority deciding truth.

Confidence would come primarily from external evidence: a failing test, a runtime trace, a fuzz
case, a static finding, a formal counterexample, an independent reproduction. Several language
models agreeing is not evidence, and the design would treat it as none. The lifecycle then joins the
loop the repository already runs: the colony finds suspicion; the system attempts reproduction; a
reproducible failure becomes a test; the fix becomes an invariant; the former bug becomes a permanent
new sense. The reason to keep it outside Polaris is the reason for the second witness: another
observer, not another subsystem that inherits the assumptions it is meant to check. The repository's
history contains a caution about this idea: at \code{v9.55} an earlier swarm apparatus that lived
inside the tree and asserted claims about itself was deleted as scaffolding. Myrmex would be held to
the opposite discipline, outside and adversarial, or not built.

\subsection{Engineering items the roadmap already names}

A hash-based signer (SLH-DSA) behind the same custody interface, for a rotation away from lattice
assumptions; the evaluation of a newer proof toolkit against the shipped one, gated on its
stabilising and on the epoch pipeline justifying a rewrite; the epoch pipeline at scale with
incremental Merkle maintenance and parallel proving; status distribution at content-delivery scale;
multi-region disaster recovery; the mobile-licence and verifiable-credential
format bridges, explicitly as formats and not trust models; the physical token and enrolment kit;
and the external gates that are not the repository's to supply: a penetration test, a cryptographic
audit, pilots with consenting users, certification, and institutionally independent witnesses.
